\documentclass[preview]{standalone}

\usepackage{amsmath}
\usepackage{amssymb}
\usepackage{stellar}
\usepackage{definitions}

\begin{document}

\id{psicologia-della-salute}
\genpage

\section{Psicologia della salute}

\begin{snippetdefinition}{psicologia-salute-definition}{Psicologia della salute}
    La \emph{psicologia della salute} è un ambito di studio e di ricerca della
    psicologia finalizzato a contribuire alla comprensione, alla promozione e al
    mantenimento della salute psicologica e fisica.
\end{snippetdefinition}

\begin{snippet}{matarazzo-psicologia-salute}
    Secondo Matarazzo, la psicologia della salute comprende contributi
    scientifici, professionali e formativi alla promozione e al mantenimento
    della salute, alla prevenzione e al trattamento della malattia e
    all'identificazione dei correlati eziologici e diagnostici della salute,
    della malattia e delle disfunzioni associate.
\end{snippet}

\begin{snippet}{nascita-psicologia-salute}
    La psicologia della salute nasce anche in risposta al cambiamento delle cause
    principali di mortalità, ai fattori economici della cura e alla necessità di
    integrare l'approccio medico tradizionale con il contributo della psicologia.
\end{snippet}

\begin{snippet}{settori-psicologia-salute}
    I principali settori della psicologia della salute riguardano promozione e
    mantenimento della salute, prevenzione e trattamento della malattia,
    correlati eziologici e diagnostici della salute e della malattia, sistema di
    cura e politiche sanitarie.
\end{snippet}

\begin{snippetexample}{adolescenti-fumo-example}{Adolescenti e fumo}
    Per studiare i fattori psicosociali che inducono gli adolescenti a fumare,
    gli studi descrittivi descrivono relazioni tra variabili, come stato di
    salute e comportamenti funzionali o disfunzionali. Gli studi predittivi
    cercano invece di prevedere nel tempo la relazione tra comportamenti e salute.
    Gli interventi traducono le conoscenze teoriche in strategie per prevenire
    comportamenti dannosi.
\end{snippetexample}

\begin{snippetdefinition}{modello-biomedico-definition}{Modello biomedico}
    Il \emph{modello biomedico} spiega la malattia esclusivamente sulla base di
    cause di natura organica, ben definite e identificabili.
\end{snippetdefinition}

\begin{snippetdefinition}{modello-biopsicosociale-definition}{Modello biopsicosociale}
    Il \emph{modello biopsicosociale} è un orientamento concettuale che sostiene
    un approccio multidimensionale alla salute e alla malattia, considerate come
    il risultato di una complessa articolazione di piani biologici, psichici e
    sociali.
\end{snippetdefinition}

\includesnpt[width=62\%,src=/snippet/static-psychology/biopsychosocial-health-model.jpg]{centered-img}

\begin{snippet}{modello-biopsicosociale-sistemi}
    Il modello biopsicosociale risente della teoria generale dei sistemi: il
    sistema è un'unità composta da parti in relazione, le parti tendono
    all'equilibrio e un cambiamento in una parte influenza la totalità del
    sistema. In questa prospettiva, la gravità della malattia dipende anche dal
    vissuto soggettivo e dal processo interpretativo dell'individuo.
\end{snippet}

\begin{snippetdefinition}{medicina-psicosomatica-definition}{Medicina psicosomatica}
    La \emph{medicina psicosomatica} studia i rapporti tra funzionamento della
    psiche e funzionamento del corpo in un'ottica terapeutica.
\end{snippetdefinition}

\begin{snippet}{medicina-psicosomatica-note}
    La medicina psicosomatica presuppone una visione olistica dell'uomo e del
    paziente. Considera variabili psicologiche, sociali e biologiche e riconosce
    che conflitti psichici e meccanismi emozionali possono incidere sulla salute
    fisica.
\end{snippet}

\begin{snippetdefinition}{health-literacy-definition}{Health literacy}
    La \emph{health literacy} è l'insieme delle conoscenze di base di area
    medica attraverso cui l'individuo mette in atto comportamenti virtuosi nella
    gestione del proprio benessere psicofisico.
\end{snippetdefinition}

\begin{snippet}{health-literacy-rischi}
    Bassi livelli di health literacy possono peggiorare lo stato di salute,
    aumentare disuguaglianze e costi sanitari e favorire comportamenti dannosi,
    accessi impropri ai servizi di emergenza o credenze magico-fatalistiche
    sulla salute.
\end{snippet}

\end{document}
